\chapter{Introduction}
\label{chapter:introduction}

\section{Single-photon detection with superconducting nanowires}
\label{sec:intro-snspd}

\emph{Why single photons matter (quantum communication, imaging, ...); what an SNSPD is and why it leads in
efficiency, jitter, dark counts and count rate \cite{esmaeilzadeh2021,natarajan2012,venza2025}.}

\section{The trade-off between efficiency and dark counts}
\label{sec:intro-tradeoff}

\emph{Internal detection efficiency against bias; the plateau; dark counts rising towards the switching current
\cite{engel2015}. The operating window in normalised bias $i_b = I_b/I_c$ as the figure of merit of this thesis.}

\section{Engineered defects: holes, slots and constrictions}
\label{sec:intro-defects}

\emph{Current crowding and geometry-dependent critical currents \cite{clem2011}; hole-patterned microstrips
\cite{li2022,li2025}; pinhole grids and their TDGL description \cite{zhu2023,zhu2025}; what is missing for
nanowire widths comparable to the hole size and to $\xi$. Why the project moved from hole arrays to slot
columns.}

\section{Research question and approach}
\label{sec:intro-question}

\emph{The research question and sub-questions; the approach (generalised TDGL with a heat model and a readout
circuit, solved with pyTDGL \cite{bishopvanhorn2023}); what is and is not claimed.}

\section{Outline}
\label{sec:intro-outline}

\emph{One sentence per chapter: \cref{chapter:theory} theory, \cref{chapter:methods} the numerical model,
\cref{chapter:results} results, \cref{chapter:experiment} experiment, \cref{chapter:discussion} discussion,
\cref{chapter:conclusion} conclusion.}
